% Proposed B177, separate late38 consolidation. No shared or frozen source edit.
\subsection*{The complete path-four contact image and its actual compact refinement}
\statement{B.177}{Proved}{The entire five-generator contact image in the two-square quotient is a path-four group with full ambient profinite topology, and its exact marked quotient strictly refines the earlier joint source-kernel test}
\label{stmt:B177}
All profinite completions in this statement are full completions.
Use B.175's actual index-two quotient
\(\Gamma_0^*\twoheadrightarrow G_\square\), and write its presentation as
\begin{equation}\label{eq:B177-target}
\begin{gathered}
 F=F(U,k,q,y),\qquad S_*=F(m,V),\qquad K_0=F\rtimes S_*,\\
 J_*=\langle S_*,U\rangle,\qquad
 H_*=\langle S_*,kUk^{-1}\rangle,\\
 K_1=K_0*_{S_*}(S_*\times\mathbb Z_{D_0}),\qquad
 L=\langle J_*,D_0H_*D_0^{-1}\rangle,\\
 G_\square=K_1*_{L}(L\times\mathbb Z_R).
\end{gathered}
\end{equation}
These are the whole-source groups of B.175: both signed source actions
and the complete node saturation are retained. The actual contact
images are
\begin{equation}\label{eq:B177-marking}
 (A_s,m,B_s,D_s,E_s)
       \longmapsto(R,m,U^{-1},UD_0U^{-1},1).
\end{equation}
Let \(J_m=\langle m,U\rangle\), \(K_b=\langle J_m,D_0\rangle\),
and let \(W\) be the image of the actual five-generator source
\(H_5\) in \(G_\square\). Then the exact ordinary presentations are
\begin{equation}\label{eq:B177-ordinary}
\begin{gathered}
 J_m=F(m,U),\qquad
 K_b=\langle m,U,D_0\mid[m,D_0]=1\rangle,\\
 W=\langle m,U,D_0,R\mid[m,D_0]=[R,m]=[R,U]=1\rangle.
\end{gathered}
\end{equation}
Thus \(W\) is the graph group of the four-vertex path
\(U\!-
 R\!-
 m\!-
 D_0\). It has its full induced profinite topology in
the actual ordinary group \(G_\square\).

Let \(P_0=(\mathbb Z_A\times\mathbb Z_m)*F(B,D,E)\),
and let \(C_0^{\mathrm{ord}}=\ker(P_0\to H_5)\).
For its actual compact source map put
\[
 \widehat P_0=\widehat{\mathbb Z}^{\,2}_{A,m}
                 \amalg\widehat F(B,D,E),\qquad
 C_0=\ker(\widehat P_0\twoheadrightarrow\mathcal H_5),
 \qquad\mathcal H_5=\overline{H_5}\subseteq\widehat{\Gamma_0}.
\]
Here the letters \(A,B,D,E\) denote the original marked source
letters. Let \(\mathcal R_i\) be closed normal subgroups of
\(\widehat P_0\) defined by
\begin{equation}\label{eq:B177-kernels}
\begin{aligned}
 \mathcal R_1&=\overline{\langle\!\langle[m,D]\rangle\!\rangle},\\
 \mathcal R_2&=\overline{\langle\!\langle
                  A,E,[m,BDB^{-1}]\rangle\!\rangle},\\
 \mathcal R_3&=\overline{\langle\!\langle
                  E,[A,B],[m,BDB^{-1}]\rangle\!\rangle}.
\end{aligned}
\end{equation}
There is an actual marked continuous epimorphism
\(\mathcal H_5\twoheadrightarrow\widehat W\), and
\begin{equation}\label{eq:B177-compact}
 \widehat P_0/\mathcal R_3\cong\widehat W,
 \qquad C_0\subseteq\mathcal R_3,
 \qquad
 \ker(\mathcal H_5\to\widehat W)=\mathcal R_3/C_0.
\end{equation}
The epimorphism \(\widehat W\twoheadrightarrow
\widehat P_0/\mathcal R_2\) is exactly the earlier B.173 contact
map, with kernel the closed normal closure of \(A\) in \(\widehat W\).
Moreover
\begin{equation}\label{eq:B177-strict}
 C_0\subseteq\mathcal R_1\cap\mathcal R_3
       \ \subsetneq\ \mathcal R_1\cap\mathcal R_2.
\end{equation}
The strict inclusion is witnessed by B.175's already separated word
\(w=[[m,D],A]\), with \([a,b]=aba^{-1}b^{-1}\).
No equality of \(C_0\) with the remaining intersection and no
fixed-source surjectivity conclusion is asserted.

\begin{proof}
\emph{The entire ordinary intersection calculation.}
The original source action restricts on \(F(U,q)\) to
\[
 m(U)=q^{-1}Uq,\qquad m(q)=q.
\]
The map \(F(U,q)\rtimes\langle m\rangle\to F(U,q)\)
which fixes \(U,q\) and sends \(m\) to \(q^{-1}\) respects these
relations. Its restriction to the subgroup generated by \(m,U\)
sends them to the free basis \(q^{-1},U\). Thus \(J_m=F(m,U)\).
The projection to \(S_*\) gives
\begin{equation}\label{eq:B177-Nq}
 J_m\cap F=N_q
   :=\langle q^nUq^{-n}:n\in\mathbb Z\rangle,
 \qquad J_m\cap S_*=\langle m\rangle.
\end{equation}
The displayed elements are the usual Schreier basis of the kernel
of \(F(m,U)\to\langle m\rangle\), after the same substitution.

Put \(T=F(k,q,y)\) and \(N=\ker(F\to T)\), where \(U\) is
killed. The ordinary Schreier basis of \(N\) is
\[
                 U_g=gUg^{-1}\qquad(g\in T).
\]
Let \(\eta:T\to Q=(C_2)^3\) record the three exponent parities
\(k,q,y\), and put \(O=\{k,q+y\}\subset Q\).
For clarity, B.175's exact affine label actions have linear columns
\[
 M_m=(k+q,q,y+q),\qquad M_V=(q+y,q,k+q),
\]
and translations \(q\) and \(0\), respectively. The first affine
map fixes both points of \(O\), and the second exchanges them.
Both are bijections, so the inverse actions preserve \(O\) as well.
Since each source sends \(U\) to a conjugate of \(U\), every signed
source translate of \(kUk^{-1}\) is an \(N\)-conjugate of a basis
letter \(U_g\) with \(\eta(g)\in O\). This retains any
\(N\)-component of the conjugator; it does not replace an actual
source action by a descended action on a finite wreath group.

Let \(N_O\) be the normal closure \emph{inside \(N\)} of those
\(O\)-labelled basis letters. Moving all source letters to the right
in \(H_*\) proves \(H_*\cap F\subseteq N_O\).
All basis letters \(U_{q^n}\) in \(N_q\) have labels \(0\) or
\(q\), outside \(O\). The basis of \(N\) therefore splits into
the \(q^n\)-letters, the \(O\)-letters, and the remaining letters.
The free-basis quotient killing the \(O\)-letters is injective on
the free factor \(N_q\), whence \(N_O\cap N_q=1\).
For \(h\in H_*\cap J_m\), remove its source coordinate \(m^n\).
The remainder belongs to \((H_*\cap F)\cap N_q=1\). Hence
\begin{equation}\label{eq:B177-HJm}
                       H_*\cap J_m=\langle m\rangle.
\end{equation}
This is an ordinary free-basis argument, with no assertion about the
abstract completion of \(N\) or \(N_q\).

The first group \(K_1\) in \eqref{eq:B177-target} is the HNN
extension with stable letter \(D_0\) centralizing \(S_*\).
Because \(J_m\cap S_*=\langle m\rangle\), every reduced word of
\(J_m*_{\langle m\rangle}
(\langle m\rangle\times\mathbb Z_{D_0})\) remains Britton-reduced
in \(K_1\). This proves the claimed presentation of \(K_b\),
and \(K_b\cap K_0=J_m\).
Likewise a reduced alternating word proves the actual embedding
\begin{equation}\label{eq:B177-Lamalgam}
                L=J_* *_{S_*}D_0H_*D_0^{-1}.
\end{equation}
Indeed its nontrivial alternating factors produce a Britton-reduced
word with internal \(H_*\setminus S_*\) and
\(J_*\setminus S_*\) coefficients.

We claim \(L\cap K_b=J_m\). An element of \(L\setminus J_*\)
has a reduced form
\[
 j_0D_0h_1D_0^{-1}j_1\cdots D_0h_rD_0^{-1}j_r,
 \qquad h_i\in H_*\setminus S_*,
\]
with internal \(j_i\in J_*\setminus S_*\).
A reduced form in \(K_b\) has coefficients in \(J_m\), and each
coefficient between opposite stable-letter signs lies outside
\(\langle m\rangle\). Equality of two reduced HNN forms forces
their stable-letter signs to agree and corresponding coefficients
to differ by adjacent edge-group multipliers. To see this directly,
apply Britton's lemma to one form times the inverse of the other and
remove the forced pinch successively. Thus at a
\(D_0h_iD_0^{-1}\) segment, the corresponding \(K_b\) coefficient
would belong to \(S_*h_iS_*\). Since \(S_*\le H_*\), it would
belong to \(H_*\cap J_m=\langle m\rangle\), contradicting
reducedness. No such segment is possible. The intersection element
lies in \(J_*\) and in \(K_b\cap K_0=J_m\). Consequently
\begin{equation}\label{eq:B177-LKb-ordinary}
                           L\cap K_b=J_m.
\end{equation}
The second HNN extension now embeds
\(K_b*_{J_m}(J_m\times\mathbb Z_R)\) by the same Britton
argument. Its image is exactly the group generated by the contacts
in \eqref{eq:B177-marking}, since these recover \(U,D_0,R,m\).
This proves all ordinary presentations in \eqref{eq:B177-ordinary}.

\emph{Continuous free families and actual split kernels.}
For a profinite space \(X\), let \(\mathfrak F(X)\) denote the
profinite group with a continuous generator family indexed by \(X\)
and the universal property for continuous maps from \(X\) into
profinite groups. It is constructed as the inverse limit of the
finite-rank free profinite groups on finite clopen partitions of
\(X\), with folding maps. The continuous family generates every
coordinate and hence topologically generates the limit. A continuous
family into a finite group factors through one finite partition;
this gives the unique homomorphism. Passing through the compatible
finite quotients of an arbitrary profinite target proves the full
universal property.

We will use its split-kernel form. For a profinite group \(T_c\),
the actual kernel of the retraction
\(\mathfrak F(X)\amalg T_c\to T_c\) is
\begin{equation}\label{eq:B177-free-kernel}
 \mathfrak F(X\times T_c),\qquad
 e_{x,g}\longmapsto g u_xg^{-1}.
\end{equation}
The \(T_c\)-action is left translation on the second index.
Finite quotients of \(T_c\) and finite partitions of \(X\) give
compatible models with finite indexing sets, proving joint
continuity. For the two inverse presentation maps, send \(u_x\)
to \(e_{x,1}\) and fix \(T_c\) in one direction; send
\(e_{x,g}\) to \(gu_xg^{-1}\) and fix \(T_c\) in the other.
Equivariance makes these homomorphisms of the semidirect products,
and both composites fix every continuous generator. This proves
\eqref{eq:B177-free-kernel} inside the actual compact group,
without completing an infinite-index discrete kernel.

\emph{The actual completed first centralizer extension.}
Characteristic finite-index subgroups of the finitely generated
free group \(F\) are cofinal. They supply actual finite semidirect
quotients of \(K_0\), so its fiber has full topology and its source
action extends continuously. The two presentation maps yield
\[
                       \widehat K_0=\widehat F\rtimes\widehat S_*.
\]
For \(n\in\widehat{\mathbb Z}\), write
\(U_n,k_n,q_n,y_n\) for four continuous families, and set
\[
 \mathcal F_D=
 \mathfrak F(\widehat{\mathbb Z}\times\{U,k,q,y\}).
\]
Let \(\widehat D_0=\widehat{\mathbb Z}\) translate \(n\), and
let \(S_*\) act on every four-letter factor by its original action.
These actions commute and extend jointly continuously. Here is the
cofinality needed for that assertion. Fold \(n\) modulo any positive
integer to get a finitely generated free profinite group with the
simultaneous source action and a commuting cyclic permutation.
Every continuous finite target of the four generator families factors
through one such folding, by local constancy with a common modulus.
In that finite-rank free profinite group its kernel contains an open
characteristic subgroup, obtained by intersecting all open subgroups
of a sufficiently bounded index. Such refinements are cofinal and
invariant under both actions. They prove the asserted continuity
without assuming that an infinite-rank automorphism group is compact.

The exact compact presentation is
\begin{equation}\label{eq:B177-K1-completion}
 \widehat K_1=
       \mathcal F_D\rtimes(\widehat S_*\times\widehat D_0).
\end{equation}
The forward map sends the original four fiber generators to their
\(n=0\) copies and sends the two section factors identically.
The reverse map sends \(z_n\) to \(D_0^n z D_0^{-n}\) and fixes
the source and stable-letter factors. These are continuous families.
The commuting relation and the actual source action prove
equivariance; both composites fix all generators and families.
Thus \eqref{eq:B177-K1-completion} is a direct model of the
completion of the displayed ordinary presentation.

Let \(\mathcal T_D\) be the quotient of \(\mathcal F_D\) killing
all \(U_n\); it is the continuous free group on the three head
families. Its actual split kernel, by
\eqref{eq:B177-free-kernel}, is
\begin{equation}\label{eq:B177-ND}
 \mathcal N_D=\ker(\mathcal F_D\to\mathcal T_D)
       =\mathfrak F(\widehat{\mathbb Z}\times\mathcal T_D),
 \qquad e_{n,g}=gU_ng^{-1}.
\end{equation}
In particular the head index ranges over the \emph{whole}
\(\mathcal T_D\), not only one four-letter factor.

\emph{The full compact edge intersection.}
In \(\widehat F\) let \(\mathcal N_J,\mathcal N_H\) be the
closed subgroups generated by all actual source translates of
\(U\) and \(kUk^{-1}\), respectively. They lie in the actual
kernel obtained by killing \(U\), and
\[
 \overline{J_*}=\mathcal N_J\rtimes\widehat S_*,\qquad
 \overline{H_*}=\mathcal N_H\rtimes\widehat S_*.
\]
These are actual closures: each compact product contains the ordinary
group densely, and its two factors are contained in that closure.
Moving source letters to the right in \(L\), using
\([D_0,S_*]=1\), similarly gives in
\eqref{eq:B177-K1-completion}
\begin{equation}\label{eq:B177-L-closure}
 \overline L=\mathcal N_L\rtimes\widehat S_*,\qquad
 \mathcal N_L=
       \overline{\langle\mathcal N_{J,0},\mathcal N_{H,1}\rangle}.
\end{equation}
The numeric subscripts are the \(D_0\)-indices.
The closure of the actual \(K_b\) is
\begin{equation}\label{eq:B177-Kb-closure}
 \begin{gathered}
 \overline K_b=
   \mathcal N_b\rtimes
       (\widehat{\mathbb Z}_m\times\widehat D_0),\\
 \mathcal N_b=
 \overline{\langle q_n^{-z}U_nq_n^z:
                       n,z\in\widehat{\mathbb Z}\rangle}.
 \end{gathered}
\end{equation}
Indeed the ordinary kernel after killing \(U\) is generated by
the integer \(m,D_0\)-conjugates of \(U\). Their actual formulas
are the displayed words. Their closed generated subgroup is
normalized by the compact section, and its product with that
section is precisely the closure by density and compactness.
The section is faithful because it projects identically to its
factor of \(\widehat S_*\times\widehat D_0\).

Define \(\eta:\mathcal T_D\to(C_2)^3\) by sending every
\(k_n,q_n,y_n\) to \(k,q,y\), respectively. On the continuous
family in \eqref{eq:B177-ND} define
\begin{equation}\label{eq:B177-erasure}
 r(e_{n,g})=
 \begin{cases}
  1,& n\equiv1\pmod2\text{ and }\eta(g)\in O,\\
  e_{n,g},&\text{otherwise}.
 \end{cases}
\end{equation}
The rule is continuous since both cases are clopen, and therefore
defines a continuous endomorphism of the actual \(\mathcal N_D\).
It fixes \(\mathcal N_b\), because the generators of
\eqref{eq:B177-Kb-closure} have head parity \(0\) or \(q\).
It fixes \(\mathcal N_{J,0}\), since it fixes all family letters
at index zero. It kills \(\mathcal N_{H,1}\): the entire signed
orbit calculation above writes every ordinary orbit generator as a
conjugate of a killed \(O\)-labelled letter, and continuity then
kills its closure. The conjugating factors themselves are retained;
no continuous choice of such factors is needed. Consequently
\[
 r|_{\mathcal N_b}=\mathrm{id},\qquad
                    r(\mathcal N_L)=\mathcal N_{J,0}.
\]

An element of \(\overline L\cap\overline K_b\) projects into
\(\widehat{\mathbb Z}_m\times\{0\}\). Removing this actual
section leaves \(z\in\mathcal N_L\cap\mathcal N_b\), so
\(z=r(z)\in\mathcal N_{J,0}\).
Now fold every even-index four-letter factor of \(\mathcal F_D\)
identically onto the factor at zero, and send every odd-index
factor to one. This is another continuous family map. It fixes the
zero factor, and sends each generator of \(\mathcal N_b\) either
to \(q^{-z}Uq^z\) or to one. Identifying the zero factor with
\(\widehat F\), the element just obtained therefore belongs to
\[
 \mathcal N_q=
     \overline{\langle q^{-z}Uq^z:z\in\widehat{\mathbb Z}\rangle}.
\]
The converse inclusion, including the \(m\)-section, is immediate.
We have proved the exact intersection in the actual completion:
\begin{equation}\label{eq:B177-LKb-compact}
                   \overline L\cap\overline K_b=\overline J_m.
\end{equation}
Neither erasure map was asserted to extend to the whole ambient group.

\emph{Full topology on the initial subgroup \(J_m\).}
For an arbitrary finite homomorphism \(\psi:J_m\to A_f\), put
\(\mu=\psi(m)\), \(r_0=\operatorname{ord}(\mu)\), and define
\[
 \eta_\psi:F\to A_f,
 \qquad U\mapsto\psi(U),\quad q\mapsto\mu^{-1},\quad k,y\mapsto1.
\]
Choose a characteristic finite-index subgroup \(C\) of \(F\)
contained in \(\ker\eta_\psi\). Retain both the finite
\(S_*\)-image in \(\operatorname{Aut}(F/C)\) and the source
character to \(C_{r_0}\) sending \(m\) to one and \(V\) to zero.
Their combined image gives an actual finite semidirect quotient
of \(K_0\). For \(g\in J_m\), its semidirect expression is
\(g=fm^n\) with \(f\in F(U,q)\), and the actual formulas for
\(m\) give
\[
                          \psi(g)=\eta_\psi(f)\mu^n.
\]
This follows on generators and their products, because
\(\eta_\psi(m(U))=\mu\eta_\psi(U)\mu^{-1}\) and the same
identity holds for \(q\). If \(g\) dies in the finite quotient,
then \(f\in C\) and \(n\equiv0\pmod{r_0}\), hence
\(\psi(g)=1\). Every finite \(J_m\)-kernel is therefore refined
by an actual finite \(K_0\)-kernel. This proves full induced
topology, without a retraction \(K_0\to J_m\) or an assertion
that \(\eta_\psi\) is \(S_*\)-equivariant on all of \(F\).
In the actual \(\widehat K_0\), the dense-product argument gives
\[
 \overline J_m=\mathcal N_q\rtimes\widehat{\mathbb Z}_m,
 \qquad
 \overline J_m\cap\widehat S_*=\widehat{\mathbb Z}_m.
\]

\emph{Finite-target extension across a centralizer extension.}
We give the complete extension lemma used twice. Let \(A_0\)
be an ordinary group and \(B_0,H_0\) actual subgroups, with
\(C_0'=B_0\cap H_0\). Suppose \(B_0\) has full induced
topology in \(A_0\), that
\(\overline{B_0}\cap\overline{H_0}=\overline{C_0'}\) in
\(\widehat A_0\), and that the ordinary subgroup generated by
\(B_0,t\) in \(A_0*_{H_0}(H_0\times\mathbb Z_t)\) is
the embedded group \(B_0*_{C_0'}(C_0'\times\mathbb Z_t)\).
Then this subgroup has full induced topology.

To prove it, take any finite homomorphism \(\psi\) from that
subgroup into \(F_0\). Its restriction to \(B_0\) extends
continuously to \(\overline{B_0}\). There is an open normal
subgroup \(M\) of \(\widehat A_0\) such that
\begin{equation}\label{eq:B177-finite-intersection}
 \begin{gathered}
 M\cap\overline{B_0}\subseteq\ker(\psi|_{B_0}),\\
 \overline{B_0}\cap\overline{H_0}M
      \subseteq\psi^{-1}(\psi(C_0')).
 \end{gathered}
\end{equation}
The first follows from the full topology. For the second, the
compact sets on the left decrease to
\(\overline{B_0}\cap\overline{H_0}=\overline{C_0'}\),
contained in the indicated clopen subset of \(\overline{B_0}\).
Compactness gives one \(M\), which can also be refined for the
first condition.

Let \(A_f=A_0/(A_0\cap M)\) and let \(B_f,H_f\) be its
subgroup images. The first condition makes \(\psi|_{B_0}\)
factor through \(B_f\); the second says
\(\psi(B_f\cap H_f)\subseteq\psi(C_0')\).
On the finite induced set
\[
 \Omega=A_f\times_{B_f}F_0,
 \qquad (ab,z)\sim(a,\psi(b)z),
\]
let \(A_f\) act on the left. On the \(H_f\)-orbit of the
identity fiber set
\[
                         t([h,z])=[h,\psi(t)z],
\]
and let \(t\) be identity on the complement. An ambiguity in
\(h\) belongs to \(H_f\cap B_f\); its \(\psi\)-image lies
in \(\psi(C_0')\) and commutes with \(\psi(t)\). Thus the
formula is well-defined. Its inverse uses \(\psi(t)^{-1}\),
and it commutes with \(H_f\) on the orbit and on its invariant
complement. This is an actual finite permutation target of the
whole centralizer extension. The identity fiber is invariant under
\(B_0,t\), and their restrictions are the left regular actions of
their \(\psi\)-images. Its kernel on the specified subgroup is
therefore contained in \(\ker\psi\). This proves the lemma;
it refines the kernel and does not assert extension to the same
finite codomain \(F_0\).

Apply the lemma first to
\[
 (A_0,B_0,H_0,C_0',t)
       =(K_0,J_m,S_*,\langle m\rangle,D_0).
\]
The ordinary embedding, full \(J_m\)-topology and actual compact
intersection were all proved above. Hence \(K_b\) has full
topology in \(K_1\). Apply it again with
\[
 (A_0,B_0,H_0,C_0',t)=(K_1,K_b,L,J_m,R).
\]
The ordinary intersection is \eqref{eq:B177-LKb-ordinary}; the
compact intersection is \eqref{eq:B177-LKb-compact}. It follows
that \(W\) has full induced topology in \(G_\square\).

\emph{The actual source quotient and strict refinement.}
Every finite \(W\)-kernel is now refined by a finite
\(G_\square\)-kernel. Pull back through the actual map of
\(\Gamma_0^*\), and take its finite-index core in the actual
index-two overgroup \(\Gamma_0\). This preserves the refinement
on \(H_5\). Equivalently the completed quotient maps its actual
compact source onto the closure of \(W\), which is
\(\widehat W\) by the full-topology proof. No equality between
\(\widehat{H_5}\) and \(\mathcal H_5\) is needed.

Substituting \(U=B^{-1}\), \(R=A\), \(D_0=BDB^{-1}\)
in \eqref{eq:B177-ordinary} gives the exact marked presentation
\[
 W=\langle A,m,B,D\mid[A,m]=[A,B]=[m,BDB^{-1}]=1\rangle,
 \qquad E\mapsto1.
\]
The substitutions and their inverses fix every marked generator.
Its identical full profinite universal presentation therefore
gives \(\widehat P_0/\mathcal R_3\cong\widehat W\).
The actual factorization through \(\mathcal H_5\) proves
\eqref{eq:B177-compact}; this is an equality of actual closed
kernels, not just a list of vanishing words. Writing \(R_3^{\mathrm{ord}}\) for the ordinary normal closure
of \(E,[A,B],[m,BDB^{-1}]\) in \(P_0\), the analogous ordinary
statement is
\[
 P_0/R_3^{\mathrm{ord}}\cong W,\qquad
 C_0^{\mathrm{ord}}\subseteq R_3^{\mathrm{ord}},\qquad
 \ker(H_5\to W)=R_3^{\mathrm{ord}}/C_0^{\mathrm{ord}}.
\]

Since \([A,B]\) belongs to the closed normal closure of \(A\),
\(\mathcal R_3\subseteq\mathcal R_2\); conversely
\(\mathcal R_2\) is exactly the closed normal closure of
\(\mathcal R_3\) and \(A\). Killing \(A\) in the last
presentation leaves exactly B.173's original marked target.
Both composite maps from \(\mathcal H_5\) agree on its five
topological generators, hence everywhere. The quotient kernel is
the closed normal closure of \(A\) in \(\widehat W\).
The older actual first target gives \(C_0\subseteq\mathcal R_1\).
Finally \(w=[[m,D],A]\) lies in both \(\mathcal R_1\) and
\(\mathcal R_2\), whereas B.175's actual finite separator makes
its image in \(\widehat W\) nonidentity. Thus \(w\notin
\mathcal R_3\), proving \eqref{eq:B177-strict}.
\end{proof}

\begingroup\footnotesize\raggedright
This consolidates the complete independently reviewed proofs in
\path{work_CE/ordinary_path_four_contact_image38.txt},
\path{work_DE/compact_path_four_contact_image38.md}, and
\path{work_CE/path_four_refines_second_contact_target38.txt}.
Their exact source bindings and the proposed-source hash are retained
in the accompanying B177 manifest. No new mathematical producer is
used: the only computational predecessor is B.175's already executed
515-check certificate. The strengthened source-image presentation
does not prove B.164's cofinal contact-image containment or the separate
original B.75 global zero-set condition. Both main questions remain open.
\par\endgroup
